\chapter{Testing}

\section{Introduction}

Testing is a critical phase in the SafeSurf system to ensure the phishing detection model is accurate, reliable, and efficient. Because the system uses a hybrid machine learning approach, DNS-based filtering, and real-time API processing, I tested each component thoroughly.

\noindent
The testing phase focuses on:
\begin{itemize}
    \item Evaluating model performance.
    \item Validating real-time detection capabilities.
    \item Ensuring correct handling of edge cases.
    \item Verifying system stability under varying network conditions.
\end{itemize}

\section{Dataset Analysis}

\subsection*{Comparative Analysis of Datasets}
\noindent
During development, multiple URL datasets were analyzed, including:
\begin{enumerate}
    \item \textbf{PhishTank Dataset:} Contains real-world, active phishing URLs \cite{ref4}.
    \item \textbf{PhiUSIIL Phishing URL Dataset:} A balanced, comprehensive dataset containing over 235,000 verified phishing and legitimate URLs \cite{ref5}.
\end{enumerate}

\subsection*{Final Decision}
\noindent
The system uses a combined dataset approach. Extracting exactly 111 features across a large, diverse dataset improves:
\begin{itemize}
    \item Model generalization.
    \item Detection accuracy across new, unseen URLs.
    \item Overall system stability.
\end{itemize}
\noindent
Testing showed that smaller datasets resulted in over-fitting and poor detection of modern, complex phishing patterns.

\section{Three-Phase Retraining and Validation Workflow}
To ensure the system remains resilient against evolving phishing tactics, a standardized three-phase pipeline was established for model retraining and performance verification. This workflow automates the transition from raw URL data to a production-ready ensemble model.

\subsection*{Phase 1: High-Concurrency Feature Extraction}
\noindent
Using the \texttt{generate\_training\_dataset.py} utility, the system processes large-scale URL repositories (e.g., the PhiUSIIL 200k URL dataset). 
\begin{itemize}
    \item \textbf{Parallelism}: Utilizes a 40-thread parallel processing pool to minimize extraction latency across hundreds of thousands of rows.
    \item \textbf{Domain-Level Caching}: Implements an in-memory cache to prevent redundant infrastructure lookups for identical domains within the same dataset.
    \item \textbf{Checkpointing}: Automatically saves progress every 5,000 processed rows to prevent data loss during long-running extraction tasks.
\end{itemize}

\subsection*{Phase 2: Multi-Model Ensemble Training}
\noindent
The \texttt{train\_deep\_clean.py} script executes a rigorous machine learning pipeline to generate the final predictive binary.
\begin{itemize}
    \item \textbf{Data Partitioning}: Implements a stratified split (70\% training, 10\% calibration, and 20\% independent holdout testing).
    \item \textbf{Model Fusion}: Deploys a \texttt{VotingClassifier} integrating Logistic Regression, Linear SVC, Random Forest, HistGradientBoosting, and XGBoost.
    \item \textbf{Class Balancing}: Uses \texttt{scale\_pos\_weight} to mathematically adjust for imbalances between legitimate and phishing samples in the source data.
\end{itemize}

\subsection*{Phase 3: Production Sanity Validation}
\noindent
Before deployment, the \texttt{validate\_clean\_model.py} utility performs a Safe-Domain Stress Test to ensure zero-tolerance for false positives on high-authority infrastructure.
\begin{itemize}
    \item \textbf{Whitelist Coverage}: Tests the trained model against a rotating list of global authorities including \url{https://google.com}, \url{https://wikipedia.org}, and \url{https://github.com}.
    \item \textbf{Assertion Logic}: The validation script triggers a build failure if any high-reputation domain is classified as phishing, guaranteeing that safe browsing remains uninterrupted.
\end{itemize}

\section{Model Performance Observations}

\subsection*{Key Findings During Testing}

\begin{enumerate}
    \item \textbf{Model Variability:}
    Different machine learning models (like Logistic Regression vs. Random Forest) showed varying performance depending on the URL structure. The ensemble approach stabilized these differences.

    \item \textbf{Feature Importance Extraction:}
    During testing, specific features proved to be highly significant for classification accuracy. These included:
    \begin{itemize}
        \item Lexical lengths and the presence of special characters.
        \item Domain age and WHOIS registration timing.
        \item SSL certificate validation.
    \end{itemize}

    \item \textbf{Impact of Ensemble Learning:}
    Single models occasionally produced false positives. Using the \textbf{VotingClassifier ensemble} (fused with an NLP model at a 65\% / 35\% weight split) helped:
    \begin{itemize}
        \item Reduce false positives.
        \item Improve mathematical reliability.
    \end{itemize}

    \item \textbf{Effect of Real-World Data:}
    Testing against live URLs revealed that attackers constantly change their URL patterns. This confirmed the need for continuous model retraining and drift monitoring.
\end{enumerate}

\section{Test Case Identification and Results}
\noindent
The following execution scenarios validate the system's performance against diverse real-world inputs, ranging from trusted global authorities to sophisticated phishing attempts and infrastructure-level anomalies. Each test case reflects a specific threat vector or legitimate browsing pattern encountered by the Chrome Extension.

\begin{itemize}
    \item \textbf{Test Case TC-01: Global Authority Validation}
    \begin{itemize}
        \item \textbf{Scenario:} Identifying high-reputation global domains.
        \item \textbf{Input URL:} \url{https://www.google.com}
        \item \textbf{Verdict:} Safe (Low Risk)
        \item \textbf{Result:} Successfully classified with a near-zero risk score due to long registration history and valid SSL certificates.
    \end{itemize}

    \item \textbf{Test Case TC-02: Region-Specific Authority}
    \begin{itemize}
        \item \textbf{Scenario:} Validating localized e-commerce portals.
        \item \textbf{Input URL:} \url{https://www.amazon.in/}
        \item \textbf{Verdict:} Safe (Low Risk)
        \item \textbf{Result:} Successfully validated using localized DNS records and legitimate infrastructure signatures.
    \end{itemize}

    \item \textbf{Test Case TC-03: Deep Directory Navigation}
    \begin{itemize}
        \item \textbf{Scenario:} Testing lexical resilience against long, complex directory paths.
        \item \textbf{Input URL:} \url{https://github.com/user/repo/tree/main/src/main.py}
        \item \textbf{Verdict:} Safe (Low Risk)
        \item \textbf{Result:} The model correctly differentiated between a complex file path on a trusted domain and a deceptive phishing directory.
    \end{itemize}

    \item \textbf{Test Case TC-04: Typo Squatting Detection}
    \begin{itemize}
        \item \textbf{Scenario:} Detecting brands spoofed via character substitution (e.g., '0' for 'o').
        \item \textbf{Input URL:} \url{g00gle.com}
        \item \textbf{Verdict:} Invalid (NXDOMAIN)
        \item \textbf{Result:} The DNS Guard intercepted the request immediately, identifying that the character-spoofed domain did not resolve to a valid IP.
    \end{itemize}

    \item \textbf{Test Case TC-05: Expired or Unregistered Domain}
    \begin{itemize}
        \item \textbf{Scenario:} Handling dormant or disposable threat domains.
        \item \textbf{Input URL:} \url{secure-login-alert.xyz}
        \item \textbf{Verdict:} Invalid (NXDOMAIN)
        \item \textbf{Result:} Correctness confirmed; the system prevents unnecessary ML execution on dead infrastructure.
    \end{itemize}

    \item \textbf{Test Case TC-06: User-info Obfuscation Attempt}
    \begin{itemize}
        \item \textbf{Scenario:} Detecting URLs using '@' to hide the actual malicious host.
        \item \textbf{Input URL:} \url{bank-login@verify.com}
        \item \textbf{Verdict:} Invalid (Invalid Host)
        \item \textbf{Result:} The Lexical Layer identified the malformed user-info pattern and flagged the host as suspicious.
    \end{itemize}

    \item \textbf{Test Case TC-07: Brand Sandwiching (Phishing)}
    \begin{itemize}
        \item \textbf{Scenario:} Detecting brand names embedded in malicious subdomains.
        \item \textbf{Input URL:} \url{google.com.login-verify.top}
        \item \textbf{Verdict:} Phishing (HIGH Risk)
        \item \textbf{Result:} The NLP pipeline and XGBoost ensemble successfully identified the deceptive use of "google.com" in the subdomain prefix.
    \end{itemize}

    \item \textbf{Test Case TC-08: IP-Based Hosting (Phishing)}
    \begin{itemize}
        \item \textbf{Scenario:} Detecting phishing sites hosted directly on IPv4 addresses.
        \item \textbf{Input URL:} \url{http://185.199.108.153/login}
        \item \textbf{Verdict:} Phishing (HIGH Risk)
        \item \textbf{Result:} The Infrastructure Layer immediately flagged the lack of a valid domain name and associated SSL certificate.
    \end{itemize}
\end{itemize}

\noindent
\textbf{Summary of Testing Results:}
\noindent
As an additional summary resource, the following table provides a condensed view of the execution statuses for the above cases:

\begin{table}[H]
\centering
\footnotesize
\renewcommand{\arraystretch}{1.5}
\caption{System Validation Summary Table}
\begin{tabularx}{\textwidth}{|l|p{3.5cm}|X|l|l|l|}
\hline
\textbf{ID} & \textbf{Scenario} & \textbf{Input URL} & \textbf{Verdict} & \textbf{Risk} & \textbf{Status} \\ \hline
TC-01 & Global Authority    & \url{google.com}       & Safe     & LOW          & \textbf{\textcolor{namegreen}{Pass}} \\ \hline
TC-02 & Region Specific     & \url{amazon.in}       & Safe     & LOW          & \textbf{\textcolor{namegreen}{Pass}} \\ \hline
TC-03 & Deep Path           & \url{github.com/...} & Safe     & LOW          & \textbf{\textcolor{namegreen}{Pass}} \\ \hline
TC-04 & Typo Squatting      & \url{g00gle.com}                   & Invalid  & NXDOMAIN     & \textbf{\textcolor{namegreen}{Pass}} \\ \hline
TC-05 & Expired Domain      & \url{...alert.xyz}       & Invalid  & NXDOMAIN     & \textbf{\textcolor{namegreen}{Pass}} \\ \hline
TC-06 & Obfuscation         & \url{bank@verify.com}       & Invalid  & Invalid Host & \textbf{\textcolor{namegreen}{Pass}} \\ \hline
TC-07 & Brand Sandwiching   & \url{google.com...}  & Phishing & HIGH         & \textbf{\textcolor{namegreen}{Pass}} \\ \hline
TC-08 & IP Hosting          & \url{185.199.108...} & Phishing & HIGH         & \textbf{\textcolor{namegreen}{Pass}} \\ \hline
\end{tabularx}
\end{table}

\noindent
\textbf{Final Testing Observations:}
\begin{itemize}
    \item \textbf{Legitimate Traffic (TC-01 to TC-03):} High-authority domains with standard SSL certificates consistently resolve as "Safe" with negligible risk scores.
    \item \textbf{Infrastructure Filtering (TC-04 to TC-06):} The DNS Guard and Lexical Normalizer successfully intercept malformed or non-existent domains before they hit the ML pipeline, reporting "Invalid" or "NXDOMAIN" status.
    \item \textbf{Heuristic Detection (TC-07 to TC-08):} Sophisticated phishing patterns were correctly flagged by the ensemble model with "HIGH" risk.
\end{itemize}

\section{Impact of System Components on Stability}

\begin{itemize}
    \item \textbf{DNS Guard Layer:}
    Instantly filters invalid or dead domains. This saves API resources and reduces unnecessary machine learning computation.

    \item \textbf{Parallel Feature Extraction:}
    Extracting 111 network features can be slow. Using \texttt{asyncio.gather} and strict 15-second timeouts prevented external query delays from freezing the API.

    \item \textbf{Ensemble Model Behavior:}
    Combining 5 distinct infrastructure algorithms with a separate NLP model drastically reduces the dependency on a single point of failure.

    \item \textbf{Threshold-Based Decision Logic:}
    Mapping raw probabilities into discrete bands (High/Medium/Low) helps developers maintain a strict balance between false positives and false negatives.
\end{itemize}

\section{Final Insights}
\noindent
Based on the completed test phases, the following conclusions were drawn:

\begin{enumerate}
    \item \textbf{Ensemble Models Improve Reliability:} Combining multiple models logically yields more stable and accurate predictions than any standalone script.
    \item \textbf{Pre-Filtering Improves Efficiency:} The deterministic DNS Guard acts as a vital shield, protecting the ML pipeline from processing junk traffic.
    \item \textbf{Feature Extraction is Critical:} The raw quality of the 111 extracted features directly dictates the eventual model performance.
    \item \textbf{Real-World Testing is Essential:} Phishing strategies evolve rapidly, requiring a system built for zero-downtime hot-reloads and continuous updates.
\end{enumerate}
